% CENG467 Final Report — CodeEnhancer Model Comparison
% LNCS format (Springer Lecture Notes in Computer Science)
\documentclass[runningheads]{llncs}

% ─── Packages ───
\usepackage[utf8]{inputenc}
\usepackage[T1]{fontenc}
\usepackage[english]{babel}
\usepackage{graphicx}
\usepackage{booktabs}
\usepackage{multirow}
\usepackage{hyperref}
\usepackage{listings}
\usepackage{xcolor}
\usepackage{float}
\usepackage{caption}
\usepackage{subcaption}
\usepackage{amsmath}
\usepackage{enumitem}

% ─── Code listing style ───
\lstset{
  language=Python,
  basicstyle=\ttfamily\footnotesize,
  keywordstyle=\color{blue}\bfseries,
  commentstyle=\color{gray},
  stringstyle=\color{red!70!black},
  frame=single,
  breaklines=true,
  numbers=left,
  numberstyle=\tiny\color{gray},
  captionpos=b
}

\hypersetup{
  colorlinks=true,
  linkcolor=blue!70!black,
  citecolor=blue!70!black,
  urlcolor=blue!70!black
}

\begin{document}

% ─── Title ───
\title{CodeEnhancer Model Comparison: Evaluating Security-Aware Code Generation Across Open-Source LLMs and Prompt Engineering Strategies}

\titlerunning{CodeEnhancer Model Comparison}

\author{Ali Ekin Özçetin \and Ali Sezer Uysal \and Mert Güner}

\authorrunning{Özçetin et al.}

\institute{Department of Computer Engineering, Izmir Institute of Technology, Izmir, Turkey \\
\small Advisor: Prof.\ Dr.\ Aytu\u{g} Onan}

\maketitle

% ============================================
\begin{abstract}
The rapid adoption of large language models (LLMs) for automated code generation has raised serious concerns about the security and reliability of LLM-produced code. This paper presents a systematic comparative study built on the CodeEnhancer framework, evaluating five open-source LLMs---Qwen 2.5 Coder 7B, DeepSeek-Coder 6.7B, Llama 3.1 8B, Gemma 2 9B, and Mistral 7B---across three prompt engineering strategies (Zero-shot, Few-shot, and Chain-of-Thought) using 81 Python prompts drawn from the LLMSecEval benchmark, yielding 1,215 experimental scenarios in total. To eliminate the self-assessment bias inherent in using a generator model as its own judge, we integrate GPT-4o-mini as an independent external LLM-Judge while keeping all local models responsible for code repair. Our results show that Qwen 2.5 Coder 7B achieves the highest average resolution rate (80.2\%), followed by DeepSeek-Coder 6.7B (77.3\%) and Llama 3.1 8B (75.7\%). Few-shot prompting with secure coding examples outperforms both Zero-shot and Chain-of-Thought across all models, and few-shot examples generalize to unseen CWE categories without data leakage. Chain-of-Thought prompting paradoxically degrades performance for smaller, general-purpose models due to context distraction under iterative repair.
\keywords{LLM security \and code generation \and prompt engineering \and SAST \and CWE \and vulnerability detection \and open-source LLMs}
\end{abstract}

% ============================================
\section{Introduction}
\label{sec:intro}

The rapid advancement of large language models such as GPT-4~\cite{openai2023gpt4}, Code Llama~\cite{roziere2024codellama}, and Codex~\cite{chen2021codex} is transforming software development. Through tools like GitHub Copilot, LLMs are integrated into daily programming workflows, accelerating development cycles and automating repetitive coding tasks. This offers substantial potential for boosting developer productivity. However, this transformative potential is tempered by critical concerns regarding the security and reliability of LLM-generated code.

Empirical studies have confirmed that code generated by popular LLMs frequently contains weaknesses listed in the MITRE Common Weakness Enumeration (CWE) Top 25~\cite{pearce2022asleep,perry2023users}. These weaknesses include command injection (CWE-78), SQL injection (CWE-89), path traversal (CWE-22), and insecure use of cryptography (CWE-327). Moreover, LLMs tend to replicate insecure coding patterns from their training data, producing code susceptible to exploits such as command injection, path traversal, and improper cryptographic function handling. Beyond security vulnerabilities, ensuring that generated code precisely implements the functionality specified in the natural language prompt---a property often referred to as functional correctness---remains a notable hurdle.

To comprehensively address these multifaceted issues, Lee et al.\ proposed CodeEnhancer~\cite{lee2026codeenhancer}, a two-stage framework that integrates LLMs with SAST tools through an iterative validation pipeline and a subsequent fine-tuning methodology. Their work demonstrated that the iterative generate-validate-refine loop, when applied to GPT-4o, eliminated 82.8\% of initial vulnerabilities on the LLMSecEval benchmark. However, CodeEnhancer evaluated only GPT-4o and relied on that same model as both generator and judge---a setup that suffers from self-assessment bias, as Olausson et al.~\cite{olausson2024selfrepair} demonstrated that LLMs frequently fail to detect or correct their own errors in the absence of external feedback.

This paper extends the CodeEnhancer framework by (1) replacing the proprietary GPT-4o generator with five open-source, locally-runnable LLMs spanning code-specialized and general-purpose architectures; (2) systematically comparing three prompt engineering strategies---Zero-shot, Few-shot, and Chain-of-Thought; and (3) decoupling the judge from the generator by using GPT-4o-mini as an independent external evaluator. All code generation and repair runs on-device via Ollama, enabling fully local, cost-free, and reproducible experimentation. Our study aims to answer: which open-source model and which prompt strategy combination produces the most secure and functionally correct Python code under an iterative SAST-based repair loop?

The remainder of this paper is organized as follows. Section~\ref{sec:related} reviews related work. Section~\ref{sec:methodology} describes our research methodology. Section~\ref{sec:experiments} presents the experimental setup and results. Section~\ref{sec:conclusion} concludes and discusses future directions.

% ============================================
\section{Related Work}
\label{sec:related}

\subsection{LLM-Based Code Generation and Security Vulnerabilities}

The security implications of LLM-generated code have been studied extensively since the release of GitHub Copilot. Pearce et al.~\cite{pearce2022asleep} conducted a systematic evaluation across 89 scenarios covering MITRE's Top 25 CWEs and found that approximately 40\% of the 1,689 generated programs contained security vulnerabilities. A follow-up study by Pearce et al.~\cite{pearce2022zeroshot} investigated the zero-shot vulnerability repair capabilities of Codex and found that while the model could fix some vulnerabilities, it introduced new ones in many cases. Perry et al.~\cite{perry2023users} provided a user study confirming that developers using AI code assistants write significantly less secure code and are less likely to identify vulnerabilities in their own code. Tambon et al.~\cite{tambon2024bugs} further examined bugs in LLM-generated code empirically, documenting that security flaws remain pervasive across multiple model generations.

More recently, Hajipour et al.~\cite{hajipour2024codelmsec} introduced the CodeLMSec benchmark to systematically evaluate security vulnerabilities in black-box code language models, covering 13 CWEs across 280 insecure prompts. Dai et al.~\cite{dai2025comprehensive} conducted a comprehensive survey of LLM secure code generation, documenting patterns of vulnerability replication across model families.

\subsection{Prompt Engineering for Code Generation}

Prompt engineering has emerged as a key lever for improving LLM performance without fine-tuning. Brown et al.~\cite{brown2020gpt3} introduced the few-shot in-context learning paradigm in GPT-3, demonstrating that prepending a small number of input-output examples to the prompt produces strong task performance across diverse NLP tasks. Wei et al.~\cite{wei2022cot} proposed Chain-of-Thought prompting, instructing models to produce intermediate reasoning steps before arriving at a final answer, which substantially improves multi-step reasoning. However, Wei et al.\ also observed that CoT gains were predominantly seen in models exceeding 100B parameters, suggesting that smaller models may not benefit as reliably.

For code generation specifically, several studies have explored how prompt structure affects output quality. Li et al.~\cite{li2024finetuning} demonstrated that fine-tuning on security patch commits yields measurable improvements in vulnerability reduction. Weyssow et al.~\cite{weyssow2025peft} explored parameter-efficient fine-tuning (PEFT) methods for code generation, focusing more on syntactic and structural quality. In the security domain, Alrashedy et al.~\cite{alrashedy2024canlllms} evaluated whether LLMs can patch security issues through prompted guidance, finding that structured prompting substantially improves vulnerability mitigation success rates compared to bare zero-shot requests.

\subsection{SAST-Based Iterative Repair}

Static Application Security Testing (SAST) tools have been integrated with LLMs to create iterative repair loops. Bandit~\cite{pycqa2014bandit} is a Python-specific SAST tool designed to detect security vulnerabilities through AST-based analysis, targeting patterns such as use of unsafe functions, hardcoded credentials, and SQL injection sinks. Pylint~\cite{pycqa2003pylint} addresses syntactic and stylistic correctness. Keltek et al.~\cite{keltek2024lsast} proposed LSAST, a hybrid system combining SAST tools with LLMs to enhance vulnerability detection coverage. Alrashedy et al.~\cite{alrashedy2024canlllms} also demonstrated that SAST-driven prompting---feeding Bandit output directly into the LLM repair prompt---is one of the earliest and most effective pipelines to combine static analysis and generation in an automated feedback loop.

The CodeEnhancer framework~\cite{lee2026codeenhancer} formalized this approach into a structured generate-validate-refine loop, achieving 82.8\% vulnerability resolution on LLMSecEval with GPT-4o. Our work extends this loop to open-source local models with varying prompt strategies.

\subsection{LLM Self-Assessment Bias}

A critical limitation in prior iterative repair systems is the use of the same LLM as both generator and judge. Olausson et al.~\cite{olausson2024selfrepair} demonstrated in a rigorous ICLR 2024 study that self-repair gains are marginal and inconsistent for code generation: when accounting for the cost of repair attempts, performance improvements are modest and sometimes absent, because models are fundamentally limited in their ability to produce accurate feedback about their own outputs. Huang et al.~\cite{huang2023large} further documented that large language models cannot reliably self-correct their reasoning without external feedback. Kamoi et al.~\cite{kamoi2024when} provided a critical survey of LLM self-correction, concluding that reliable self-correction only occurs when expert feedback---not autonomous introspection---is provided.

This motivates our design choice to decouple the judge from the generator, using GPT-4o-mini as an independent external evaluator. Our mid-phase results validate this decision empirically: Llama 3.1 8B achieved 0\% resolution when judging its own outputs, but recovered to 75.7\% average resolution after the judge was decoupled.

\subsection{Evaluation Benchmarks}

The LLMSecEval benchmark~\cite{tony2023llmseceval} provides 150 natural language prompts covering 18 of MITRE's CWE Top 25, each designed to elicit a specific vulnerability type when processed by an LLM. SecurityEval~\cite{siddiq2022securityeval} offers a broader set of 130 samples across 75 vulnerability types, useful for generalization testing. Our study uses the 81 Python-language prompts from LLMSecEval, filtering out C-language prompts incompatible with our Python-specific SAST pipeline.

% ============================================
\section{Research Methodology}
\label{sec:methodology}

\subsection{Baseline Framework: CodeEnhancer Stage 1}

Our study builds directly on Stage 1 of the CodeEnhancer framework~\cite{lee2026codeenhancer}. The original Stage 1 pipeline consists of: (1) structured code generation using a prompt template that enforces an embedded docstring with \texttt{Input Prompt}, \texttt{Intention}, and \texttt{Functionality} fields; (2) iterative validation using Bandit (security) and Pylint (syntax) to detect issues and generate structured feedback reports; and (3) LLM-driven code refinement based on the feedback, repeated up to a maximum of 5 iterations. The original framework used GPT-4o as both generator and judge.

We extend Stage 1 in three dimensions: (i) replacing GPT-4o with five open-source local models running via Ollama; (ii) adding three distinct prompt engineering strategies at the generation step; and (iii) decoupling the functional correctness judge from the generator model by substituting GPT-4o-mini as an independent external evaluator.

\subsection{Decoupled LLM-Judge Design}

The key architectural modification in our pipeline is the separation of the generator and the judge. In the original CodeEnhancer, the same GPT-4o model that generates code also evaluates its functional correctness---a setup vulnerable to self-assessment bias~\cite{olausson2024selfrepair}. In our pipeline, the local model is responsible only for code generation and SAST-based repair. Once a generated file passes all SAST checks (no Bandit or Pylint issues), it is submitted to GPT-4o-mini for functional correctness evaluation via a structured JSON output mode. The judge evaluates whether the code correctly implements the docstring specification, returning \texttt{"verdict": "Correct"} or \texttt{"verdict": "Incorrect"} with a \texttt{"reason"} field. If the verdict is \texttt{Incorrect}, the reason is forwarded to the local generator model for a functional fix, and the file re-enters the SAST loop. This design ensures that functional correctness evaluation is fully decoupled from the generator's inherent biases.

\subsection{Prompt Engineering Strategies}

We implement three prompt strategies applied at the initial code generation step:

\paragraph{Zero-shot.} The model receives the LLMSecEval natural language prompt together with a structured docstring template and formatting instructions, but no examples or explicit security guidance. System prompt: \textit{``You are a code generation assistant.''}

\paragraph{Few-shot.} Five secure code examples are prepended to the prompt, covering CWE-78 (command injection), CWE-89 (SQL injection), CWE-502 (deserialization), CWE-434 (unrestricted file upload), and CWE-22 (path traversal). These examples demonstrate both the expected docstring structure and concrete secure coding patterns (input sanitization, parameterized queries, safe deserialization). System prompt: \textit{``You are a secure code generation assistant.''}

\paragraph{Chain-of-Thought (CoT).} The model is instructed to follow a four-step reasoning process before writing code: (1) identify all external inputs; (2) identify CWE risks associated with those inputs; (3) select appropriate mitigations; (4) write secure code implementing those mitigations. System prompt: \textit{``You are a security-aware code generation assistant.''}

\subsection{Evaluation Metrics}

For each of the 1,215 scenarios (5 models $\times$ 3 strategies $\times$ 81 prompts), we measure:
\begin{itemize}
  \item \textbf{Bandit@0}: Percentage of initially generated files containing at least one Bandit finding (initial vulnerability rate).
  \item \textbf{Resolution Rate}: Percentage of files that pass all SAST and judge checks within 5 iterations (final resolution success).
  \item \textbf{Unresolved Rate}: Percentage of files still failing after 5 iterations.
  \item \textbf{Average Iterations}: Mean number of repair iterations required to reach resolution (for resolved files).
\end{itemize}

% ============================================
\section{Experimental Study}
\label{sec:experiments}

\subsection{Setup}

All experiments were executed locally on a single workstation running Ollama v0.x with OpenAI-compatible REST API at \texttt{localhost:11434/v1}. All models used temperature = 0 (deterministic), max tokens = 2048, and context window (\texttt{num\_ctx}) = 8192 to ensure identical context limits across all models and prompt strategies. GPT-4o-mini served as the external judge, with total API cost of \$0.34 for all 1,215 scenarios. SAST tools used were Bandit v1.8.3 and Pylint v3.3.7.

\subsection{Models}

Table~\ref{tab:models} summarizes the five models evaluated. The selection covers two code-specialized models (Qwen 2.5 Coder, DeepSeek-Coder) and three general-purpose models (Llama 3.1, Gemma 2, Mistral), spanning three distinct vendor families (Alibaba, Meta/DeepSeek, Google).

\begin{table}[H]
\centering
\caption{Models evaluated in the experiment pipeline.}
\label{tab:models}
\begin{tabular}{lllr}
\toprule
\textbf{Model} & \textbf{Family} & \textbf{Specialization} & \textbf{Parameters} \\
\midrule
Qwen 2.5 Coder  & Alibaba  & Code generation  & 7B   \\
DeepSeek-Coder  & DeepSeek & Code generation  & 6.7B \\
Llama 3.1       & Meta     & General-purpose  & 8B   \\
Gemma 2         & Google   & General-purpose  & 9B   \\
Mistral 7B      & Mistral  & General-purpose  & 7B   \\
\bottomrule
\end{tabular}
\end{table}

\subsection{Dataset}

We use the 81 unique Python-language prompts from the LLMSecEval benchmark~\cite{tony2023llmseceval}, covering 12 CWE categories including CWE-78, CWE-89, CWE-22, CWE-502, and CWE-327 among others. The 67 C-language prompts were excluded because our SAST pipeline depends on Python-specific tools. Stratified sampling was used to remove duplicates while preserving CWE coverage.

\subsection{Results}

\subsubsection{Per-Scenario Results}

Table~\ref{tab:main_results} presents the full results across all 15 experiment combinations.

\begin{table}[H]
\centering
\caption{Results across all 5 models $\times$ 3 strategies $\times$ 81 prompts = 1,215 scenarios.}
\label{tab:main_results}
\small
\begin{tabular}{llrrrr}
\toprule
\textbf{Model} & \textbf{Strategy} & \textbf{Bandit@0 (\%)} & \textbf{Resolved (\%)} & \textbf{Unresolved (\%)} & \textbf{Avg Iter.} \\
\midrule
\multirow{3}{*}{Qwen 2.5 Coder 7B}
  & Zero-shot & 63.0 & 80.2 & 19.8 & 2.31 \\
  & Few-shot  & 54.3 & \textbf{82.7} & \textbf{17.3} & 2.38 \\
  & CoT       & 55.6 & 77.8 & 22.2 & 2.41 \\
\midrule
\multirow{3}{*}{DeepSeek-Coder 6.7B}
  & Zero-shot & 45.7 & 80.2 & 19.8 & \textbf{2.23} \\
  & Few-shot  & 64.2 & 75.3 & 24.7 & 2.68 \\
  & CoT       & 37.0 & 76.5 & 23.5 & 2.43 \\
\midrule
\multirow{3}{*}{Llama 3.1 8B}
  & Zero-shot & 37.0 & 70.4 & 29.6 & 2.72 \\
  & Few-shot  & 40.7 & 79.0 & 21.0 & 2.57 \\
  & CoT       & 27.2 & 77.8 & 22.2 & 2.48 \\
\midrule
\multirow{3}{*}{Gemma 2 9B}
  & Zero-shot & 34.6 & 64.2 & 35.8 & 2.99 \\
  & Few-shot  & 40.7 & 71.6 & 28.4 & 2.58 \\
  & CoT       & 29.6 & 55.6 & 44.4 & 3.22 \\
\midrule
\multirow{3}{*}{Mistral 7B}
  & Zero-shot & 59.3 & 64.2 & 35.8 & 3.37 \\
  & Few-shot  & 33.3 & 58.0 & 42.0 & 3.11 \\
  & CoT       & 60.5 & 45.7 & 54.3 & 3.83 \\
\bottomrule
\end{tabular}
\end{table}

\subsubsection{Model-Level Averages}

Table~\ref{tab:model_avg} aggregates performance across all three strategies per model.

\begin{table}[H]
\centering
\caption{Model-level averages across all three prompt strategies.}
\label{tab:model_avg}
\begin{tabular}{lrrr}
\toprule
\textbf{Model} & \textbf{Avg Bandit@0 (\%)} & \textbf{Avg Resolution (\%)} & \textbf{Avg Iterations} \\
\midrule
Qwen 2.5 Coder 7B   & 57.6 & \textbf{80.2} & \textbf{2.37} \\
DeepSeek-Coder 6.7B & 49.0 & 77.3 & 2.45 \\
Llama 3.1 8B        & 35.0 & 75.7 & 2.59 \\
Gemma 2 9B          & 35.0 & 63.8 & 2.93 \\
Mistral 7B          & 51.0 & 56.0 & 3.44 \\
\bottomrule
\end{tabular}
\end{table}

\subsubsection{Strategy-Level Averages}

Table~\ref{tab:strategy_avg} aggregates performance across all five models per strategy.

\begin{table}[H]
\centering
\caption{Strategy-level averages across all five models.}
\label{tab:strategy_avg}
\begin{tabular}{lrrr}
\toprule
\textbf{Strategy} & \textbf{Avg Bandit@0 (\%)} & \textbf{Avg Resolution (\%)} & \textbf{Avg Iterations} \\
\midrule
Zero-shot        & 47.9 & 71.8 & 2.72 \\
Few-shot         & 46.6 & \textbf{73.3} & \textbf{2.66} \\
Chain-of-Thought & 42.0 & 66.7 & 2.87 \\
\bottomrule
\end{tabular}
\end{table}

\subsection{Analysis}

\subsubsection{Code-Specialized vs.\ General-Purpose Models}

Figure~\ref{fig:sec_res} shows the initial vulnerability rate (Bandit@0) and the final resolution rate side-by-side. Qwen 2.5 Coder 7B and DeepSeek-Coder 6.7B (code-specialized) achieved the highest overall resolution rates (80.2\% and 77.3\% average), while Mistral 7B lagged significantly at 56.0\%. The code-specialized models showed both higher initial resolution capability and faster convergence (lower average iterations). Notably, Llama 3.1 8B (general-purpose) performed surprisingly well at 75.7\% average resolution---a result that validates the utility of the decoupled judge, since Llama achieved 0\% resolution in mid-phase experiments when it served as its own judge.

\begin{figure}[H]
  \centering
  \begin{subfigure}[b]{0.48\textwidth}
    \centering
    \includegraphics[width=\textwidth]{figures/fig1_bandit_hit_rate.png}
    \caption{Initial vulnerability rate (Bandit@0).}
    \label{fig:bandit}
  \end{subfigure}
  \hfill
  \begin{subfigure}[b]{0.48\textwidth}
    \centering
    \includegraphics[width=\textwidth]{figures/fig2_resolution_rate.png}
    \caption{Resolution success rate.}
    \label{fig:resolution}
  \end{subfigure}
  \caption{Initial security and final resolution success comparison across all five models.}
  \label{fig:sec_res}
\end{figure}

\subsubsection{Repair Efficiency}

Figure~\ref{fig:efficiency} shows average iterations and iteration distribution. Qwen 2.5 Coder 7B is the most efficient model (2.37 average iterations), while Mistral 7B requires 3.44 average iterations. The box plot reveals that code-specialized models have lower medians and narrower spreads, converging to secure, correct code in fewer attempts. General-purpose models, especially Mistral, exhibit larger spreads and frequently reach the 5-iteration limit.

\begin{figure}[H]
  \centering
  \begin{subfigure}[b]{0.48\textwidth}
    \centering
    \includegraphics[width=\textwidth]{figures/fig3_avg_iterations.png}
    \caption{Average iterations per model (lower is better).}
    \label{fig:avg_iter}
  \end{subfigure}
  \hfill
  \begin{subfigure}[b]{0.48\textwidth}
    \centering
    \includegraphics[width=\textwidth]{figures/fig7_box_iterations.png}
    \caption{Iteration distribution box plot.}
    \label{fig:box_iter}
  \end{subfigure}
  \caption{Evaluation efficiency and iteration distribution per model.}
  \label{fig:efficiency}
\end{figure}

\subsubsection{Few-Shot Generalization to Unseen CWEs}

A key concern with few-shot prompting is data leakage: if the five examples cover the same CWEs as the test prompts, apparent improvements may reflect memorization rather than generalization. Figure~\ref{fig:spec_gen} disaggregates few-shot performance into seen CWEs (those covered by examples) and unseen CWEs (those not in the examples). Across all models, performance on unseen CWEs equals or exceeds performance on seen CWEs (e.g., Qwen 84.4\% unseen vs.\ 80.6\% seen; Llama 82.2\% unseen vs.\ 75.0\% seen). This confirms that few-shot examples teach a meta-rule of secure programming---input sanitization, parameterized queries, and safe deserialization---that models generalize to novel vulnerability contexts without overfitting to specific CWE patterns.

\begin{figure}[H]
  \centering
  \begin{subfigure}[b]{0.48\textwidth}
    \centering
    \includegraphics[width=\textwidth]{figures/fig8_code_vs_general.png}
    \caption{Code-specialized vs.\ general-purpose resolution.}
    \label{fig:code_vs_gen}
  \end{subfigure}
  \hfill
  \begin{subfigure}[b]{0.48\textwidth}
    \centering
    \includegraphics[width=\textwidth]{figures/fig9_seen_unseen_cwe.png}
    \caption{Few-Shot: seen vs.\ unseen CWE performance.}
    \label{fig:seen_unseen}
  \end{subfigure}
  \caption{Model specialization and few-shot generalization to unseen vulnerability types.}
  \label{fig:spec_gen}
\end{figure}

\subsubsection{The Chain-of-Thought Degradation Paradox}

Table~\ref{tab:strategy_avg} shows that Chain-of-Thought prompting yielded the lowest average resolution (66.7\%) despite producing the lowest initial vulnerability rate (42.0\% Bandit@0). This apparent paradox is explained by the interaction between CoT's verbose reasoning output and the iterative repair loop: for 7B--9B parameter models, generating multi-step reasoning text exhausts a significant portion of the available context window, leaving less capacity for code generation in subsequent repair iterations. This leads to instruction-following decay---models that successfully reason about vulnerabilities in the initial generation struggle to apply the same reasoning during repair, producing syntax errors or incomplete fixes. This finding is consistent with Wei et al.'s~\cite{wei2022cot} observation that CoT benefits are most pronounced in large models ($\geq$100B parameters).

\subsubsection{CWE-Level Vulnerability Analysis}

Figure~\ref{fig:heatmap_radar} presents the CWE-level initial vulnerability heatmap and a multi-dimensional radar chart. CWE-78 (OS Command Injection) and CWE-502 (Deserialization) are the most consistently challenging categories, triggering Bandit findings across all models and strategies. The radar chart confirms that Qwen 2.5 Coder 7B dominates in resolution efficiency and overall score, while Llama 3.1 8B excels in security rate (fewest initial Bandit flags) and unseen CWE generalization.

\begin{figure}[H]
  \centering
  \begin{subfigure}[b]{0.48\textwidth}
    \centering
    \includegraphics[width=\textwidth]{figures/fig4_cwe_heatmap.png}
    \caption{CWE-level Bandit@0 heatmap.}
    \label{fig:heatmap}
  \end{subfigure}
  \hfill
  \begin{subfigure}[b]{0.48\textwidth}
    \centering
    \includegraphics[width=\textwidth]{figures/fig6_radar_per_model.png}
    \caption{Multi-dimensional radar comparison per model.}
    \label{fig:radar}
  \end{subfigure}
  \caption{CWE vulnerability profiles and multi-dimensional performance radar.}
  \label{fig:heatmap_radar}
\end{figure}

\subsubsection{Judge Verdict Distribution}

Figure~\ref{fig:status_judge} shows the final resolution breakdown alongside the distribution of GPT-4o-mini judge verdicts (Correct vs.\ Incorrect). Mistral 7B triggered the most Incorrect verdicts across all strategies, indicating a tendency to generate functionally incomplete code even when SAST checks pass. Qwen 2.5 Coder and DeepSeek-Coder received the highest proportions of Correct verdicts, demonstrating superior docstring specification adherence.

\begin{figure}[H]
  \centering
  \begin{subfigure}[b]{0.48\textwidth}
    \centering
    \includegraphics[width=\textwidth]{figures/fig5_final_status.png}
    \caption{Resolved vs.\ unresolved final status.}
    \label{fig:status}
  \end{subfigure}
  \hfill
  \begin{subfigure}[b]{0.48\textwidth}
    \centering
    \includegraphics[width=\textwidth]{figures/fig10_judge_analysis.png}
    \caption{GPT-4o-mini judge verdict distribution.}
    \label{fig:judge_decisions}
  \end{subfigure}
  \caption{Final execution status and LLM-Judge verdict distribution per model.}
  \label{fig:status_judge}
\end{figure}

\subsection{Threats to Validity}

\textbf{SAST Coverage.} Bandit and Pylint only detect specific vulnerability patterns and syntax issues. Logic flaws and semantic vulnerabilities that do not match known patterns will not be caught by SAST, and are left to the judge. Using a broader set of SAST tools (e.g., CodeQL~\cite{github2023codeql}) could improve coverage.

\textbf{Judge Calibration.} While GPT-4o-mini is independent and structurally robust, it is still a language model and may exhibit minor false positives or negatives in functional correctness evaluation.

\textbf{Quantization Effects.} Ollama runs quantized models (typically 4-bit), which may introduce minor behavioral differences compared to full-precision weights. This is a practical necessity for local execution on consumer hardware.

\textbf{Dataset Scope.} Results are derived from 81 Python prompts across 12 CWE categories. Generalization to other programming languages, longer codebases, or real-world projects remains an open question.

% ============================================
\section{Conclusion and Future Work}
\label{sec:conclusion}

This paper presented a systematic comparative study of five open-source LLMs across three prompt engineering strategies for secure Python code generation, built on the CodeEnhancer framework. Our key findings are:

\begin{enumerate}
  \item \textbf{Decoupling the judge resolves self-assessment bias.} Llama 3.1 8B went from 0\% resolution under self-assessment to 75.7\% average resolution with an independent judge, validating the architectural importance of separating generator and evaluator roles~\cite{olausson2024selfrepair}.
  \item \textbf{Code-specialized models outperform general-purpose models.} Qwen 2.5 Coder 7B (80.2\%) and DeepSeek-Coder 6.7B (77.3\%) dominate the resolution ranking, requiring fewer repair iterations.
  \item \textbf{Few-shot prompting is the most effective strategy.} At 73.3\% average resolution with 2.66 average iterations, few-shot outperforms zero-shot (71.8\%, 2.72 iter.) and CoT (66.7\%, 2.87 iter.).
  \item \textbf{Few-shot examples generalize without data leakage.} All models achieved equal or higher resolution on unseen CWEs compared to seen CWEs under few-shot setup.
  \item \textbf{CoT degrades under repair loops for small models.} The verbose reasoning output of CoT competes with repair instructions in the context window of 7B--9B models, causing instruction-following decay during iterative fixes.
\end{enumerate}

\paragraph{Future Work.} Several directions remain open. First, extending the evaluation to the full LLMSecEval benchmark (all 150 prompts including C-language) and to SecurityEval would broaden generalizability. Second, integrating larger SAST tools such as CodeQL~\cite{github2023codeql} could reduce false negatives in vulnerability detection. Third, Stage 2 fine-tuning---excluded from the current scope---could be applied to the best-performing open-source models using framework-refined code as training data, analogous to the CodeEnhancer fine-tuning approach~\cite{lee2026codeenhancer}. Fourth, exploring retrieval-augmented prompt construction---dynamically selecting few-shot examples that match the CWE type of the target prompt---may further improve few-shot generalization.

% ============================================
\begin{thebibliography}{99}

\bibitem{lee2026codeenhancer}
J.\ Lee, K.\ Mai, N.\ Ghate, T.\ Yagyu, R.\ Beuran, Y.\ Tan:
CodeEnhancer: LLM-Generated Python Code Enhancement Through SAST Integration and Fine-Tuning.
\textit{Knowledge-Based Systems} (2026).
\url{https://doi.org/10.1016/j.knosys.2026.115925}

\bibitem{openai2023gpt4}
OpenAI:
GPT-4 Technical Report.
arXiv preprint arXiv:2303.08774 (2023).

\bibitem{roziere2024codellama}
B.\ Rozière et al.:
Code Llama: Open Foundation Models for Code.
arXiv preprint arXiv:2308.12950 (2024).

\bibitem{chen2021codex}
M.\ Chen et al.:
Evaluating Large Language Models Trained on Code.
arXiv preprint arXiv:2107.03374 (2021).

\bibitem{pearce2022asleep}
H.\ Pearce, B.\ Ahmad, B.\ Tan, B.\ Dolan-Gavitt, R.\ Karri:
Asleep at the Keyboard? Assessing the Security of GitHub Copilot's Code Contributions.
In: \textit{IEEE Symposium on Security and Privacy (SP)}, pp.\ 754--768 (2022).

\bibitem{pearce2022zeroshot}
H.\ Pearce, B.\ Tan, B.\ Ahmad, R.\ Karri, B.\ Dolan-Gavitt:
Examining Zero-Shot Vulnerability Repair with Large Language Models.
In: \textit{IEEE Symposium on Security and Privacy (SP)}, pp.\ 2339--2356 (2023).

\bibitem{perry2023users}
N.\ Perry, M.\ Srivastava, D.\ Kumar, D.\ Boneh:
Do Users Write More Insecure Code with AI Assistants?
In: \textit{Proceedings of the 2023 ACM SIGSAC Conference on Computer and Communications Security (CCS '23)}, pp.\ 2785--2799 (2023).

\bibitem{tambon2024bugs}
F.\ Tambon, A.\ Moradi-Dakhel, A.\ Nikanjam, F.\ Khomh, M.C.\ Desmarais, G.\ Antoniol:
Bugs in Large Language Models Generated Code.
\textit{Empirical Software Engineering} 30(3) (2025).

\bibitem{hajipour2024codelmsec}
H.\ Hajipour, K.\ Hassler, T.\ Holz, L.\ Schonherr, M.\ Fritz:
CodeLMSec Benchmark: Systematically Evaluating and Finding Security Vulnerabilities in Black-Box Code Language Models.
In: \textit{IEEE Conference on Secure and Trustworthy Machine Learning (SaTML)}, pp.\ 684--709 (2024).

\bibitem{dai2025comprehensive}
S.\ Dai, J.\ Xu, G.\ Tao:
A Comprehensive Survey of LLM Secure Code Generation (2025).
arXiv preprint arXiv:2503.15554.

\bibitem{brown2020gpt3}
T.\ Brown et al.:
Language Models are Few-Shot Learners.
In: \textit{Advances in Neural Information Processing Systems (NeurIPS)}, vol.\ 33 (2020).

\bibitem{wei2022cot}
J.\ Wei et al.:
Chain-of-Thought Prompting Elicits Reasoning in Large Language Models.
In: \textit{Advances in Neural Information Processing Systems (NeurIPS)}, vol.\ 35 (2022).

\bibitem{li2024finetuning}
J.\ Li, A.\ Sangalay, C.\ Cheng, Y.\ Tian, J.\ Zhang:
Fine Tuning Large Language Model for Secure Code Generation.
In: \textit{Proceedings of the 2024 IEEE/ACM FORGE '24}, pp.\ 86--90 (2024).

\bibitem{weyssow2025peft}
M.\ Weyssow, X.\ Zhou, K.\ Kim, D.\ Lo, H.\ Sahraoui:
Exploring Parameter-Efficient Fine-Tuning Techniques for Code Generation with Large Language Models.
\textit{ACM Transactions on Software Engineering and Methodology} (Jan.\ 2025).

\bibitem{alrashedy2024canlllms}
K.\ Alrashedy, A.\ Aljasser, P.\ Tahmkeer, M.\ Gombolay:
Can LLMs Patch Security Issues?
arXiv preprint arXiv:2312.00024 (2024).

\bibitem{pycqa2014bandit}
PyCQA:
Bandit --- A Security-Oriented Static Analysis Tool for Python.
\url{https://github.com/PyCQA/bandit} (2014).

\bibitem{pycqa2003pylint}
PyCQA:
Pylint --- Python Static Analysis and Linting Tool.
\url{https://github.com/pylint-dev/pylint} (2003).

\bibitem{keltek2024lsast}
M.\ Keltek, R.\ Hu, M.\ Fani Sani, Z.\ Li, P.\ Anderson:
LSAST: Enhancing Cybersecurity Through LLM-Supported SAST.
arXiv preprint arXiv:2409.15735 (2024).

\bibitem{olausson2024selfrepair}
T.X.\ Olausson, J.P.\ Inala, C.\ Wang, J.\ Gao, A.\ Solar-Lezama:
Is Self-Repair a Silver Bullet for Code Generation?
In: \textit{The Twelfth International Conference on Learning Representations (ICLR)}, Vienna, Austria (2024).

\bibitem{huang2023large}
J.\ Huang, X.\ Chen, S.\ Mishra, H.S.\ Zheng, A.W.\ Yu, X.\ Song, D.\ Zhou:
Large Language Models Cannot Self-Correct Reasoning Yet.
arXiv preprint arXiv:2310.01798 (2023).

\bibitem{kamoi2024when}
R.\ Kamoi, N.\ Zhang, N.\ Zhang, J.\ Han, R.\ Zhang:
When Can LLMs Actually Correct Their Own Mistakes? A Critical Survey of Self-Correction of LLMs.
\textit{Transactions of the Association for Computational Linguistics} 12, 1417--1440 (2024).

\bibitem{tony2023llmseceval}
C.\ Tony, M.\ Mutas, N.E.D.\ Ferreyra, R.\ Scandariato:
LLMSecEval: A Dataset of Natural Language Prompts for Security Evaluations.
In: \textit{2023 IEEE/ACM 20th International Conference on Mining Software Repositories (MSR)}, pp.\ 588--592 (2023).

\bibitem{siddiq2022securityeval}
M.L.\ Siddiq, J.C.S.\ Santos:
SecurityEval Dataset: Mining Vulnerability Examples to Evaluate Machine Learning-Based Code Generation Techniques.
In: \textit{Proceedings of the 1st International Workshop on Mining Software Repositories Applications for Privacy and Security (MSR4P\&S '22)}, pp.\ 29--33 (2022).

\bibitem{github2023codeql}
GitHub:
CodeQL --- Semantic Code Analysis Engine.
\url{https://codeql.github.com/} (2023).

\end{thebibliography}

\end{document}
